%--------------------------------------------------------------------------------
%
%
%--------------------------------------------------------------------------------
\pagebreak
\section*{LDR — Load Register Reserved}
\addcontentsline{toc}{section}{LDR — Load Register Reserved}

\begin{iPageDescEntry}{Syntax}
  \texttt{LDR RegR, SignedImmed13(RegB)}
\end{iPageDescEntry}

\begin{iPageDescEntry}{Format}
    The \texttt{LDR} instruction uses the indexed instruction format. \\

    \begin{flushleft}
    \drawInstrFormat
        {0, 1, 3, 5, 6.5, 8.5, 9.5, 16}
        {\OpCodeGrpMEM, \OpCodeLDR, RegR, 0, RegB, 3, S-IMM-13}
    \end{flushleft}
\end{iPageDescEntry}

\begin{iPageDescEntry}{Description}
    The \texttt{LDR} instruction loads a value from memory into a general-purpose register
    and marks the address as "reserved" for a subsequent conditional store.
    - Only the first step of an atomic read-modify-write pair.
    - Subsequent \texttt{STRC} (store conditional) will succeed only if the reservation is still valid.
    
    May cause a memory reference trap if the address is invalid.
\end{iPageDescEntry}

\begin{iPageDescEntry}{Operation}
    \texttt{RegR <- memLoad(memOperand());} \\
    \texttt{markAddressReserved(memOperand());}  (sets reservation for STRC)
\end{iPageDescEntry}

\begin{iPageDescEntry}{Exceptions}
    - Data Alignment Trap \\
    - Memory Access Violation Trap \\
    - Memory Protection Violation Trap \\
    - Data TLB Miss Trap
\end{iPageDescEntry}

\begin{iPageDescEntry}{Notes}
    - This instruction is part of a *load-reserved / store-conditional* pair for atomic memory updates. \\
    - Reservation may be cleared by other stores to the same address by another processor. \\
    - Often used for implementing locks or atomic counters.
\end{iPageDescEntry}
